\documentclass[12pt]{article}

\usepackage{amsmath, amssymb, amsthm}
\usepackage{geometry}
\usepackage{setspace}
\usepackage{booktabs}
\usepackage{bm}
\usepackage{bbm}
\usepackage{enumitem}

\geometry{margin=1.2in}
\doublespacing

\newtheorem{assumption}{Assumption}
\newtheorem{proposition}{Proposition}
\newtheorem{corollary}{Corollary}
\newtheorem{remark}{Remark}
\newtheorem{lemma}{Lemma}
\newtheorem{definition}{Definition}

\DeclareMathOperator{\RSS}{RSS}
\DeclareMathOperator{\E}{\mathbb{E}}
\DeclareMathOperator{\Var}{Var}
\DeclareMathOperator{\tr}{tr}
\newcommand{\Normal}{\mathcal{N}}
\newcommand{\1}{\mathbbm{1}}
\newcommand{\ytd}{\tilde{y}}
\newcommand{\Dtd}{\tilde{D}}
\newcommand{\Xtd}{\tilde{X}}

\title{Section: The Specification Issue (standalone draft)}
\author{}
\date{}

\begin{document}
\maketitle

% =============================================================================
%  SECTION: THE SPECIFICATION ISSUE
% =============================================================================
\section{The Specification Issue}
\label{sec:spec_issue}

A fundamental but underappreciated tension runs through the staggered
difference-in-differences literature.
On one side stands the fully flexible TWFE estimator, which treats
every cohort-time cell as a distinct parameter.
On the other side stands the fully pooled TWFE estimator, which collapses
all cells into a single coefficient.
Both approaches are suboptimal in general.
The fully flexible estimator is guaranteed to be unbiased but wastes
statistical precision whenever some cohort-time effects happen to be
equal in the population.
The fully pooled estimator is efficient but produces biased estimates of
every cohort-time effect the moment any two of those effects genuinely
differ.
The truth, in most empirical settings, lies somewhere between these two
extremes: a subset of cohort-time cells share a common treatment effect
and can be jointly estimated, while the remainder are heterogeneous and
must be kept separate.
This section makes that intuition precise, derives the efficiency cost
of the two extremes in closed form, and shows that any correctly specified
intermediate model strictly dominates the fully flexible estimator in terms
of estimator variance without sacrificing unbiasedness.

\bigskip

\noindent\textbf{The two extremes.}
Recall the within-transformed model obtained after double-demeaning to
partial out unit and time fixed effects:
\begin{equation}
  \ytd = \Dtd\,\bm{\tau} + \tilde{\varepsilon},
  \qquad
  \tilde{\varepsilon} \sim \bigl(\mathbf{0},\,\sigma^2 I_{NT}\bigr),
  \label{eq:flex_dm}
\end{equation}
where $\ytd$ is the $NT \times 1$ vector of within-transformed outcomes,
$\Dtd = [\tilde{d}_1, \ldots, \tilde{d}_K]$ is the $NT\times K$ matrix of
within-transformed cohort-time dummies, and $\bm{\tau} =
(\tau_1,\ldots,\tau_K)'$ collects all $K$ CATTs.
Ordinary least squares on \eqref{eq:flex_dm} yields the \emph{fully
flexible} estimator
\begin{equation}
  \hat{\bm{\tau}}_{\text{flex}}
  = \bigl(\Dtd'\Dtd\bigr)^{-1}\Dtd'\ytd,
  \qquad
  \Var\!\bigl(\hat{\bm{\tau}}_{\text{flex}}\bigr)
  = \sigma^2\bigl(\Dtd'\Dtd\bigr)^{-1}.
  \label{eq:flex_var}
\end{equation}
The diagonal element $\sigma^2[(\Dtd'\Dtd)^{-1}]_{kk}$ is the sampling
variance of the $k$-th CATT estimate.
Because each CATT is estimated from the variation in its own
within-transformed indicator $\tilde{d}_k$, this variance is governed
by how much identifying variation CATT $k$ possesses individually.
In large panels with many cohorts and long post-treatment windows, $K$
is large and the variation attributable to any single CATT cell is
correspondingly modest, leading to imprecise estimates even when the
overall sample size $NT$ is large.

The opposite extreme imposes a single common coefficient $\tau$ for every
cohort-time cell.
Writing $\tilde{X}_{\text{pool}} = \sum_{k=1}^K \tilde{d}_k$ for the
aggregate within-transformed treatment indicator, the \emph{fully pooled}
estimator is
\begin{equation}
  \hat\tau_{\text{pool}}
  = \bigl(\tilde{X}_{\text{pool}}'\tilde{X}_{\text{pool}}\bigr)^{-1}
    \tilde{X}_{\text{pool}}'\ytd,
  \qquad
  \Var\!\bigl(\hat\tau_{\text{pool}}\bigr)
  = \frac{\sigma^2}{\bigl\|\tilde{X}_{\text{pool}}\bigr\|^2}.
  \label{eq:pool_var}
\end{equation}
Since $\tilde{X}_{\text{pool}}$ aggregates the variation of all $K$ cells,
$\|\tilde{X}_{\text{pool}}\|^2 \gg \|\tilde{d}_k\|^2$ for any individual
$k$, and the pooled estimator is correspondingly precise.
The cost is bias: the pooled estimator is consistent for
$\sum_{k=1}^K w_k \tau_k^*$, the variation-weighted average of true CATTs,
not for any individual $\tau_k^*$.
Formally, the bias of $\hat\tau_{\text{pool}}$ as an estimator of $\tau_k^*$ is
\begin{equation}
  \E\bigl[\hat\tau_{\text{pool}}\bigr] - \tau_k^*
  = \sum_{j=1}^K w_j\bigl(\tau_j^* - \tau_k^*\bigr),
  \qquad
  w_j = \frac{\|\tilde{d}_j\|^2}{\sum_{\ell=1}^K \|\tilde{d}_\ell\|^2},
  \label{eq:pool_bias}
\end{equation}
where we use the approximation that the within-transformed dummies are
nearly orthogonal across distinct cohort-time cells (which holds exactly in
balanced panels).
The bias in \eqref{eq:pool_bias} is zero only when all true CATTs are equal
--- precisely the case in which the pooled model is correctly specified.
In any other case the pooled estimator distorts every cohort-specific
treatment effect toward a global average, rendering it useless for
heterogeneity-aware policy analysis.

\bigskip

\noindent\textbf{The general intermediate case.}
Let the true parameter vector $\bm{\tau}^* = (\tau_1^*,\ldots,\tau_K^*)'$
exhibit \emph{partial homogeneity}: exactly $K-p$ cells share a common
treatment effect and the remaining $p$ cells are heterogeneous, with
$0 < p < K$.
More precisely, there exists a partition $\mathcal{P}^* =
\{C_1^*,\ldots,C_p^*,\, C_{p+1}^*\}$ of $\{1,\ldots,K\}$ such that
\begin{align}
  \bigl|C_l^*\bigr| &= 1 \quad \text{for } l = 1,\ldots,p
    && \text{($p$ heterogeneous singletons)},
  \label{eq:singletons}\\
  \bigl|C_{p+1}^*\bigr| &= K - p
    && \text{($K-p$ homogeneous cells)},
  \label{eq:homo_group}\\
  \tau_k^* &= \phi_l^* \quad \text{for all } k \in C_l^*,
  \label{eq:true_effects}
\end{align}
where $\phi_1^*, \ldots, \phi_p^*$ are all distinct, $\phi_{p+1}^*$ is the
common effect shared by the homogeneous group, and all $p+1$ values are
distinct from one another.
We call $\mathcal{P}^*$ the \emph{true partition} and $K - p \geq 2$
(otherwise there is no pooling gain).
The true model has $m^* = p + 1$ distinct parameters rather than $K$,
so the fully flexible model is over-parameterised by $K - p - 1$ parameters.

To estimate the correctly specified model under $\mathcal{P}^*$, define the
group regressors $\Xtd_l = \sum_{k \in C_l^*} \tilde{d}_k$ and collect them
into the $NT \times (p+1)$ matrix
$\Xtd^* = [\tilde{X}_1^*, \ldots, \tilde{X}_p^*, \tilde{X}_{p+1}^*]$.
For the $p$ singletons, $\Xtd_l = \tilde{d}_{k_l}$.
For the homogeneous group, $\Xtd_{p+1} = \sum_{k \in C_{p+1}^*} \tilde{d}_k$.
The restricted model is
\begin{equation}
  \ytd = \Xtd^*\bm{\phi}^* + \tilde{\varepsilon},
  \qquad
  \tilde{\varepsilon} \sim \bigl(\mathbf{0},\,\sigma^2 I_{NT}\bigr),
  \label{eq:oracle_model}
\end{equation}
and its OLS estimator is
\begin{equation}
  \hat{\bm{\phi}}^* = \bigl(\Xtd^{*\prime}\Xtd^*\bigr)^{-1}\Xtd^{*\prime}\ytd,
  \qquad
  \Var\!\bigl(\hat{\bm{\phi}}^*\bigr)
  = \sigma^2\bigl(\Xtd^{*\prime}\Xtd^*\bigr)^{-1}.
  \label{eq:oracle_var}
\end{equation}
The $(p+1)$-th diagonal element of $\sigma^2(\Xtd^{*\prime}\Xtd^*)^{-1}$
governs the precision of $\hat\phi_{p+1}^*$, the estimate of the common
homogeneous effect.
For any $k \in C_{p+1}^*$, the oracle's estimate of $\tau_k^* =
\phi_{p+1}^*$ is $\hat\phi_{p+1}^*$, which pools the identifying variation
of all $K-p$ cells in the homogeneous group.
The flexible estimator, in contrast, estimates $\tau_k^*$ using only
the variation in $\tilde{d}_k$.
This difference in the effective information set is the source of the
variance gap.

\bigskip

\noindent\textbf{Variance comparison: oracle versus fully flexible.}
We now establish that the oracle estimator $\hat{\bm{\phi}}^*$ strictly
dominates the fully flexible estimator $\hat{\bm{\tau}}_{\text{flex}}$
in terms of estimator variance for every CATT in the homogeneous group.
The key insight is that both estimators are unbiased for $\phi_{p+1}^*$
under the true model, so the Gauss--Markov theorem immediately determines
which is more efficient.

\begin{lemma}[Unbiasedness of the flexible estimator under the oracle model]
  \label{lem:unbiased}
  Fix any $k \in C_{p+1}^*$.  Under the oracle DGP
  \eqref{eq:oracle_model}, the flexible OLS estimator $\hat\tau_k^{\emph{flex}}$
  is an unbiased estimator of $\phi_{p+1}^*$.
\end{lemma}

\begin{proof}
Let $e_k$ denote the $k$-th standard basis vector in $\mathbb{R}^K$, so
that $\hat\tau_k^{\text{flex}} = e_k'\hat{\bm{\tau}}_{\text{flex}} =
e_k'(\Dtd'\Dtd)^{-1}\Dtd'\ytd$.
This is a linear function of $\ytd$, so its expectation under the oracle
model equals $e_k'(\Dtd'\Dtd)^{-1}\Dtd'\,\E[\ytd]$.
Under \eqref{eq:oracle_model}, $\E[\ytd] = \Xtd^*\bm{\phi}^*$.
Observe that $\Xtd^*\bm{\phi}^*$ can be written as
$\sum_{l=1}^{p+1}\phi_l^* \Xtd_l^* = \sum_{k=1}^K \tau_k^* \tilde{d}_k =
\Dtd\bm{\tau}^*$, since every column of $\Dtd$ appears in exactly one
group of $\mathcal{P}^*$ and the corresponding group effect is precisely
$\tau_k^*$.  Therefore
\[
  \E\bigl[\hat\tau_k^{\text{flex}}\bigr]
  = e_k'\bigl(\Dtd'\Dtd\bigr)^{-1}\Dtd'\Dtd\,\bm{\tau}^*
  = e_k'\bm{\tau}^* = \tau_k^* = \phi_{p+1}^*. \qedhere
\]
\end{proof}

\noindent Armed with Lemma~\ref{lem:unbiased}, the variance ordering follows
directly from the Gauss--Markov theorem.

\begin{proposition}[Oracle efficiency dominance]
  \label{prop:oracle_dom}
  Under the true partition $\mathcal{P}^*$ and for any $k \in C_{p+1}^*$,
  \begin{equation}
    \Var\!\bigl(\hat\phi_{p+1}^*\bigr)
    \;\leq\;
    \Var\!\bigl(\hat\tau_k^{\emph{flex}}\bigr),
    \label{eq:var_ineq}
  \end{equation}
  with strict inequality whenever $|C_{p+1}^*| = K - p \geq 2$.
\end{proposition}

\begin{proof}
The oracle model \eqref{eq:oracle_model} is correctly specified.
By the Gauss--Markov theorem, $\hat\phi_{p+1}^*$ is the Best Linear
Unbiased Estimator (BLUE) of $\phi_{p+1}^*$ within the class of all
linear functions of $\ytd$.
By Lemma~\ref{lem:unbiased}, $\hat\tau_k^{\text{flex}}$ is also a linear
unbiased estimator of $\phi_{p+1}^*$.
It therefore satisfies $\Var(\hat\tau_k^{\text{flex}}) \geq
\Var(\hat\phi_{p+1}^*)$.
Strict inequality holds whenever $\hat\tau_k^{\text{flex}} \neq
\hat\phi_{p+1}^*$ with positive probability, which is the case as long as
$|C_{p+1}^*| \geq 2$, since $\hat\phi_{p+1}^*$ pools information from all
$K-p$ cells in the homogeneous group while $\hat\tau_k^{\text{flex}}$
uses only the variation in $\tilde{d}_k$.
\end{proof}

\noindent Proposition~\ref{prop:oracle_dom} applies identically to each of
the $p$ heterogeneous singleton groups: for $k \in C_l^*$ with
$l \leq p$, the oracle and flexible estimators coincide ($\hat\phi_l^* =
\hat\tau_k^{\text{flex}}$), so there is no variance gap.
The entire efficiency gain comes from pooling the $K-p$ cells in the
homogeneous group.

\bigskip

\noindent\textbf{Quantifying the efficiency gain.}
To obtain an explicit expression for the variance reduction in
\eqref{eq:var_ineq}, it is instructive to work under the additional
assumption that the within-transformed cohort-time dummies are mutually
orthogonal, $\tilde{d}_j'\tilde{d}_k = 0$ for $j \neq k$, which holds
exactly for balanced panels in which each unit belongs to exactly one
cohort.
Under orthogonality, $\Dtd'\Dtd = \mathrm{diag}(n_1,\ldots,n_K)$ where
$n_k = \|\tilde{d}_k\|^2$ denotes the effective sample size for CATT $k$.

For the flexible estimator, the variance of $\hat\tau_k^{\text{flex}}$ simplifies to
\begin{equation}
  \Var\!\bigl(\hat\tau_k^{\text{flex}}\bigr)
  = \frac{\sigma^2}{n_k},
  \qquad k = 1,\ldots,K.
  \label{eq:flex_var_orth}
\end{equation}
For the oracle estimator of the homogeneous group, orthogonality of the
individual dummies implies
\begin{equation}
  \bigl\|\tilde{X}_{p+1}^*\bigr\|^2
  = \Bigl\|\sum_{k\in C_{p+1}^*}\tilde{d}_k\Bigr\|^2
  = \sum_{k\in C_{p+1}^*} n_k,
  \label{eq:oracle_norm}
\end{equation}
so the oracle variance is
\begin{equation}
  \Var\!\bigl(\hat\phi_{p+1}^*\bigr)
  = \frac{\sigma^2}{\displaystyle\sum_{k\in C_{p+1}^*} n_k}.
  \label{eq:oracle_var_orth}
\end{equation}
Taking the ratio of \eqref{eq:oracle_var_orth} to \eqref{eq:flex_var_orth}
for any $k \in C_{p+1}^*$,
\begin{equation}
  \frac{\Var(\hat\phi_{p+1}^*)}{\Var(\hat\tau_k^{\text{flex}})}
  = \frac{n_k}{\displaystyle\sum_{j\in C_{p+1}^*} n_j}
  < 1.
  \label{eq:var_ratio}
\end{equation}
The ratio equals $n_k / (n_k + \sum_{j \neq k, j \in C_{p+1}^*} n_j)$,
which is strictly less than one and decreasing in the effective sample sizes
of the other cells in the homogeneous group.
In the balanced case, where $n_k = n$ for all $k$, the ratio collapses to
\begin{equation}
  \frac{\Var(\hat\phi_{p+1}^*)}{\Var(\hat\tau_k^{\text{flex}})}
  = \frac{1}{K - p},
  \label{eq:var_ratio_balanced}
\end{equation}
so the oracle estimator is exactly $(K-p)$ times more efficient than the
flexible estimator for every cell in the homogeneous group.
The gain grows linearly in the number of cells that are correctly pooled:
pooling together $K-p$ cells concentrates $(K-p)$ times as much effective
sample size onto a single estimated coefficient as any individual flexible
estimate.

\bigskip

\noindent\textbf{Bias of the pooled estimator under partial homogeneity.}
While the oracle strictly improves on the flexible estimator without
any bias cost, the fully pooled estimator obtains even lower variance by
imposing equality across all $K$ cells --- including the $p$ that are
genuinely heterogeneous.
Under the oracle DGP \eqref{eq:oracle_model} with orthogonal dummies, the
pooled estimator converges to
\begin{equation}
  \E\bigl[\hat\tau_{\text{pool}}\bigr]
  = \sum_{k=1}^K w_k\,\tau_k^*
  = \sum_{l=1}^p w_{k_l}\,\phi_l^*
    + \phi_{p+1}^*\!\!\!\sum_{k\in C_{p+1}^*} w_k,
  \label{eq:pool_plim}
\end{equation}
where $w_k = n_k/\sum_j n_j$.
For a heterogeneous cell $k_l \in C_l^*$, $l \leq p$, the bias is
\begin{equation}
  \E\bigl[\hat\tau_{\text{pool}}\bigr] - \phi_l^*
  = \sum_{j=1,j\neq l}^p w_{k_j}(\phi_j^* - \phi_l^*)
    + \phi_{p+1}^*\!\!\!\sum_{k\in C_{p+1}^*} w_k - w_{k_l}\phi_l^*.
  \label{eq:pool_bias_hetero}
\end{equation}
For a homogeneous cell $k \in C_{p+1}^*$, the bias is
\begin{equation}
  \E\bigl[\hat\tau_{\text{pool}}\bigr] - \phi_{p+1}^*
  = \sum_{l=1}^p w_{k_l}(\phi_l^* - \phi_{p+1}^*).
  \label{eq:pool_bias_homo}
\end{equation}
Both biases are non-zero whenever any heterogeneous cell differs from the
homogeneous group ($\phi_l^* \neq \phi_{p+1}^*$ for some $l$).
This contamination effect is a consequence of pooling across genuinely
distinct effects: the treatment signals of the heterogeneous cells
``bleed into'' the pooled estimate, and the homogeneous group's effect
is estimated as a distorted mixture of all $K$ true parameters.
Notably, the bias does not vanish as $n \to \infty$; with more data the
pooled estimator converges to the wrong limit.

\bigskip

\noindent\textbf{The specification trade-off.}
The findings above admit a unified description of the three estimators and
the trade-off they face.
Table~\ref{tab:spec_tradeoff} summarises the bias and variance of each
estimator under the partial-homogeneity DGP \eqref{eq:oracle_model}
for cells in the homogeneous group $C_{p+1}^*$ (the case where the
distinction between estimators is sharpest).

\begin{table}[ht]
  \centering
  \caption{Bias and variance under partial homogeneity for cells $k \in C_{p+1}^*$.
    Balanced-panel approximation with $n_k = n$ for all $k$.}
  \label{tab:spec_tradeoff}
  \begin{tabular}{lcc}
    \toprule
    Estimator & Bias for $\tau_k^* = \phi_{p+1}^*$ & Variance \\
    \midrule
    Fully flexible  & $0$                              & $\sigma^2 / n$ \\[4pt]
    Oracle ($\mathcal{P}^*$) & $0$               & $\sigma^2 / [(K-p)\,n]$ \\[4pt]
    Fully pooled    & $\displaystyle\sum_{l=1}^p w_{k_l}(\phi_l^* - \phi_{p+1}^*)$ & $\sigma^2 / [K\,n]$ \\
    \bottomrule
  \end{tabular}
\end{table}

\noindent The table reveals the core specification issue in staggered DiD.
Moving from the flexible to the oracle specification reduces variance by a
factor of $K-p$ at zero bias cost.
Moving further to the pooled specification achieves an additional variance
reduction (by an additional factor of $K/(K-p)$) but introduces a bias that
does not vanish with sample size.
The correct choice --- the oracle --- sits strictly between the two
extremes: it is exactly as precise as the pooled estimator among all
\emph{unbiased} estimators of $\phi_{p+1}^*$, and it is exactly as unbiased
as the flexible estimator among all estimators that pool cells in $C_{p+1}^*$.

The practical challenge is that $\mathcal{P}^*$ is unknown.
A researcher who applies the fully flexible estimator by default leaves
$(K-p-1)$ degrees of freedom on the table and may report unnecessarily
wide confidence intervals.
A researcher who collapses to a single pooled coefficient introduces
bias that cannot be detected or corrected without additional structure.
What is needed is a data-driven procedure that recovers $\mathcal{P}^*$
--- or a close approximation --- from the observed data.
The $\ell_0$-penalized greedy estimator and the Bayesian Dirichlet Process
approach developed in the following sections provide two principled answers
to this problem.

\end{document}
